\documentclass{article}\begin{document}Theorem FinalCube.fullDiagram proves the mixed base/payload equation, componentwise on every object of the actual weighted span fibre at positive length $n+1$. Quantifiers: an ordinary base category, endpoints with the existing split laws, ordinary Cat-valued presheaves in the indicated universes, an arbitrary natural transformation of payloads, and any base arrow. FinalCube.full uses the actual FibreValue.eval functors and FibreValue.evalSquare base squares for both payloads, with both endpoint corrections from EvalPayload.evalSquare. The source grid uses NativeCube.presentedSquare; the target grid uses naturality of the actual iterate transformation. The left path evaluates the presented source square with the target payload, applies the target-payload base square, and reindexes the correction at the starting base. The right inserts the correction at the restricted base, applies the restricted iterate payload functor to the source-payload base square, and applies iterate naturality. Mix.joinLaw handles the middle interchange through naturality and the stated functor equalities. The source presentation preserves comparisons by its strict inverse functors, but no native comparison is replaced by identity: only their presentation's identity whisker is removed (presentedBaseY and presentedBaseZ). NativeCube.strictCube and DecoderCube.mixed supply the interior and decoder laws, and isoJoinLaw computes the component of the exact pasteIso underlying the old evalSquare. The proof uses default heartbeat settings. This closes the fixed-positive-length mixed diagram, not the all-length telescope comparison, bundled higher coherence, the empty summand, coproduct preservation, free-monad multiplication, flattening, or strict algebra recovery. Manuscript 458ff36a, lines 929--953, distinguishes the length-one telescope unit and length-zero free-monad unit.\end{document}
